\chapter{The Agent Layer}
\label{ch:agents}

\section{Roles rather than one assistant}

A single general assistant is the wrong abstraction for autonomous work, because the
prompting, the tool set and the failure handling that make an implementer effective make a
reviewer credulous. Xyra runs a team of specialised roles with distinct mandates and distinct
tool grants.

\begin{table}[h]
\centering
\small
\begin{tabularx}{\textwidth}{l X l}
\toprule
\textbf{Role} & \textbf{Mandate} & \textbf{Write access} \\
\midrule
Planner & Decompose an objective into a dependency ordered ticket graph & None \\
Architect & Produce and defend a design, challenge another design & None \\
Implementer & Complete exactly one ticket, leave tests passing & Worktree \\
Reviewer & Cross examine a diff through one lens, return a structured verdict & None \\
Verifier & Execute tests, render pages, run quality sweeps, report facts & Sandbox only \\
Repairer & Fix a specific named failure inside a snapshot & Snapshot \\
Replanner & Split a repeatedly failing ticket into smaller tickets & None \\
Scribe & Summarise, journal, produce release notes and reports & Documents \\
\bottomrule
\end{tabularx}
\caption{Agent roles. Write access is a property of the role, not of the model.}
\label{tab:roles}
\end{table}

Visual verification, quality sweeps and cross repository analysis are capabilities of the
Verifier role rather than separate roles, which keeps the role set closed at eight. The Agent
Theatre in Chapter~\ref{ch:interface} renders exactly one lane per role in this table.

The vendor behind each role is configurable, and the default deliberately places the reviewer
on a different vendor from the implementer. A model reviewing its own output shares its own
blind spots.

\section{The mission supervisor}

The supervisor is the component that makes unattended work possible. It owns a durable ticket
graph and drives it to completion without asking the developer anything it can resolve
itself.

\begin{figure}[h]
\centering
\begin{tikzpicture}[
  st/.style={draw=xyraline,fill=xyrapanel,rounded corners=2pt,minimum height=8mm,minimum width=26mm,font=\scriptsize,align=center},
  dec/.style={draw=xyralime,fill=white,diamond,aspect=2.2,inner sep=1pt,font=\scriptsize,align=center},
  arr/.style={-{Stealth[length=1.8mm]},xyragrey,font=\scriptsize}
]
\node[st] (pick) {Pick ready\\ticket};
\node[st,right=10mm of pick] (run) {Fresh agent\\session};
\node[dec,right=10mm of run] (ver) {Verified};
\node[st,right=12mm of ver] (commit) {Commit\\ticket};
\node[dec,below=9mm of ver] (att) {Attempts\\left};
\node[st,left=12mm of att] (swap) {Swap vendor,\\retry};
\node[dec,below=9mm of att] (split) {Splittable};
\node[st,left=12mm of split] (children) {Replan into\\children};
\node[st,below=9mm of split] (quar) {Quarantine,\\continue};

\draw[arr] (pick) -- (run);
\draw[arr] (run) -- (ver);
\draw[arr] (ver) -- node[above] {yes} (commit);
\draw[arr] (ver) -- node[right] {no} (att);
\draw[arr] (att) -- node[above] {yes} (swap);
\draw[arr] (att) -- node[right] {no} (split);
\draw[arr] (split) -- node[above] {yes} (children);
\draw[arr] (split) -- node[right] {no} (quar);
\draw[arr] (swap) |- (pick);
\draw[arr] (children) |- (pick);
\draw[arr] (commit) |- ++(0,10mm) -| (pick);
\draw[arr] (quar) -| ++(-52mm,0) |- (pick);
\end{tikzpicture}
\caption{Supervisor loop. Failure is a policy with four escalations, never a stop.}
\end{figure}

\subsection{Why fresh sessions}

The central design decision is that each ticket runs in a session that has never seen any
other ticket. This is counterintuitive until the arithmetic is done. A session that
accumulates context degrades in two ways: it fills the window, and it accumulates the
residue of earlier reasoning, including earlier mistakes. A fresh session with a precise
brief and a compressed structural context reaches higher first attempt success than a long
running session at ticket fifty, and the total work a mission can perform stops being bounded
by any window at all.

\subsection{Durability}

Mission state is written atomically after every state transition. The consequences are
concrete: closing the laptop pauses a mission, a crash loses at most the ticket in flight,
and resuming is a single command that reads the graph and continues. The state file is the
mission; the process is disposable.

\subsection{Budgets and brakes}

Autonomy without limits is a liability. Every mission carries hard caps on total sessions,
wall clock duration, attempts per ticket and consecutive failures, plus a stop flag the
developer can set at any moment which halts the loop after the ticket in flight completes.
Exceeding any cap halts the mission with a stated reason rather than continuing on hope.

\section{Verification}

\subsection{Isolated execution}

Verification never runs in the developer's working tree. The supervisor snapshots the
uncommitted change set into a throwaway worktree and executes the project's own test or build
command there. The developer's tree is untouched, which means verification can run while the
developer keeps working, and a failed attempt cannot leave debris behind.

\subsection{Visual verification}

For any change with a reachable page, a verifier renders it headless, captures a screenshot
and asks a vision capable model whether the result matches the stated intent. The verdict is
structured, with issues ranked by severity, which is what makes it usable inside an automatic
fix loop rather than as decoration. A design export can be supplied, in which case the
comparison is against the intended design rather than against a description.

\subsection{Adversarial review}

The council runs a diff past a rival vendor through independent lenses, by default
correctness, security and house conventions, executed in parallel. Each lens returns findings
with severity, and a policy decides whether the aggregate blocks. Policy is per repository and
per path, so code that touches money or credentials can demand more lenses and block on lower
severities than code that touches a marketing page.

\subsection{Quality assurance sweeps}

A verifier drives the running application as a hostile user: crawling same origin pages,
filling every input with type aware adversarial values, clicking aggressively, and collecting
console errors, uncaught exceptions, server errors, failed requests, accessibility problems
and load timings. Findings are deduplicated and ranked, then fed back as tickets rather than
delivered as a raw log.

\section{Skills}

Skills are the mechanism by which a team encodes its own standards so that every agent
session inherits them without being told.

\begin{table}[h]
\centering
\small
\begin{tabularx}{\textwidth}{l X}
\toprule
\textbf{Skill} & \textbf{Encodes} \\
\midrule
Autonomy & Prove before presenting: verify, look at the result, analyse impact before touching shared contracts \\
Terminal discipline & Never run an interactive command or a process that does not exit, per tool flags tabulated \\
Conventions & House rules for code, commits and interface copy \\
Review & The checklist a reviewer applies, including domain specific traps \\
Domain safety & Rules for money arithmetic, credentials, and any high consequence domain \\
Shipping & The gate before a change leaves the machine \\
Prompt discipline & Turning a vague instruction into a scoped one before acting \\
\bottomrule
\end{tabularx}
\caption{The default skill set. Teams add their own.}
\end{table}

Skills resolve in a precedence chain: workspace, then user, then bundled defaults. A workspace
skill overrides a user skill of the same name, which lets a team enforce a standard on a
repository without touching anyone's personal configuration. The Context Inspector shows
exactly which skills were injected into any given session, because an invisible instruction
that changes behaviour is a debugging nightmare.

\section{Fleet operations}

A contract rarely lives in one repository. Fleet raises every structural capability from the
repository to the group.

\begin{itemize}
  \item Repositories are registered from the developer's version control provider, not from a
  hand written file, and missing ones are cloned on registration.
  \item Cross repository search resolves a symbol or endpoint name across every member.
  \item Cross repository impact separates the repository that defines a contract from those
  that consume it, and orders the resulting work accordingly.
  \item A fleet refactor executes the same intent through every affected repository in
  dependency order, verifying each one independently and preparing linked changes.
\end{itemize}

\section{Tool surface exposed to agents}

Agents act through a bounded tool set. Every tool is auditable, every tool is capability
gated, and no tool grants ambient authority.

\begin{longtable}{@{}p{0.24\textwidth} p{0.68\textwidth}@{}}
\toprule
\textbf{Tool} & \textbf{Contract} \\
\midrule
\endhead
read\_file & Returns a file at a requested SCP tier or in full, by handle \\
search\_code & Structural search first, similarity fallback clearly labelled \\
impact & Deterministic closure for a symbol or file, by transitive distance \\
graph & Neighbourhood of a file: imports, importers, defined symbols \\
edit & Applies a transaction to a buffer, always reversible \\
run\_tests & Executes the project's test command inside a snapshot \\
verify\_patch & Full verification of the current change set, returns a verdict \\
render\_check & Renders a page and returns a structured visual verdict \\
compare\_design & Renders a page and compares it against a design export \\
qa\_sweep & Drives the running application adversarially, returns ranked defects \\
fleet\_search & Cross repository search \\
fleet\_impact & Cross repository impact analysis \\
journal & Records a decision with its rationale against the current commit \\
ask\_developer & Escalates a genuine decision as a structured choice with costs \\
\bottomrule
\caption{Agent tool surface. Anything not on this list is not reachable.}
\end{longtable}

\section{Escalation to the developer}

The supervisor escalates only when a decision genuinely belongs to the human: an
architectural fork with real trade offs, an ambiguity in the objective that changes the
outcome, or an action with consequences outside the repository. Escalation is never a
terminal question. It is a structured choice: each option carries its measured cost, its
dependency weight and a sketch of the resulting structure, and the mission continues on
unblocked tickets while the choice is pending.
